% --- Archon physics formalization source begin ---
% archon:physics
% archon:covers PhyXMiniProblems/problem_phyx_mini_0296.lean
% archon:source-report reports/phyx_mini/problem_phyx_mini_0296.source.json

\chapter{Physics problem phyx\_mini\_0296}
\label{ch:PhyXMiniProblems_problem_phyx_mini_0296}

\paragraph{Problem source.}
\#\# Physical scenario

For a particular transverse standing wave on a long string, one of the antinodes is at \$x = 0\$ and an adjacent node is at \$x = 0.10\$ m. The displacement \$y(t)\$ of the string particle at \$x = 0\$ is shown in figure, where the scale of the \$y\$ axis is set by \$y\_s = 4.0\$ cm.

\#\# Question

What is the transverse velocity of the string particle at \$x = 0.20\$ m at \$t = 1.0\$ s?

\#\# Answer choices

- A: A: 0.12 m/s

- B: B: -0.05 m/s

- C: C: -0.18 m/s

- D: D: -0.13 m/s

\#\# Figure caption (auxiliary; use the image as primary evidence)

The image is a graph depicting a sine wave plotted in a two-dimensional Cartesian coordinate system. 

- **Axes:**
  - The horizontal axis is labeled as "t (s)" which represents time in seconds.
  - The vertical axis is labeled as "y (cm)" indicating the vertical displacement in centimeters.

- **Scales and Labels:**
  - The horizontal axis is marked with intervals at 0.5, 1.0, 1.5, and 2.0 seconds.
  - The vertical axis is labeled with \textbackslash{}( y\_s \textbackslash{}) and \textbackslash{}(-y\_s \textbackslash{}), representing the maximum and minimum values of the wave.

- **Sine Wave:**
  - The wave starts at the origin (0,0) and oscillates between positive and negative values.
  - The wave has smooth, periodic curves reaching its maximum at around 1.0 seconds and its minimum around 0.5 seconds before returning to zero.

- **Grid:**
  - The graph background includes a grid, aiding in the visualization of the wave's periodicity and amplitude.

The overall theme illustrates the periodic nature of the sine function in relation to time, depicted in centimeters over a two-second interval.

\#\# Dataset metadata

Wave Properties; Numerical Reasoning, Physical Model Grounding Reasoning

\paragraph{Recorded answer/context.}
D: D: -0.13 m/s

\paragraph{Figure/image path.}
/root/proposal\_for\_physic/science-mango/phyx\_mini\_run/phyx\_data/test\_image/296.png

\paragraph{Formalization target.}
create a compiling Lean file with sorry bodies at `PhyXMiniProblems/problem\_phyx\_mini\_0296.lean`. The Lean declarations must preserve the physical quantities, dimensions or dimensional roles, figure labels, governing-law hypotheses, and final relation expressed by this problem.
Use Mathlib/Physlib names found through LeanExplore where available. If a physics API is missing, introduce faithful local abstractions rather than scalar placeholder aliases.

\begin{theorem}[Physics formalization target]
\label{thm:physics:phyx_mini_0296:target}
\lean{PhyXMiniProblems.ProblemPhyXMini0296.problem_phyx_mini_0296}
\uses{def:physics:phyx-mini-0296:phyxminiproblems-problemphyxmini0296-transversevelocityreadout, def:physics:phyx-mini-0296:phyxminiproblems-problemphyxmini0296-standingstringwavesetup, def:physics:phyx-mini-0296:phyxminiproblems-problemphyxmini0296-matchesstandingwaveproblemdata, def:physics:phyx-mini-0296:phyxminiproblems-problemphyxmini0296-matchesprimarydisplacementgraph, def:physics:phyx-mini-0296:phyxminiproblems-problemphyxmini0296-hasphysicalstandingwaveparameters, def:physics:phyx-mini-0296:phyxminiproblems-problemphyxmini0296-satisfiessinusoidalstandingwavelaws, def:physics:phyx-mini-0296:phyxminiproblems-problemphyxmini0296-answerchoice, def:physics:phyx-mini-0296:phyxminiproblems-problemphyxmini0296-displayedvelocityinmeterspersecond, def:physics:phyx-mini-0296:phyxminiproblems-problemphyxmini0296-recordeddatasetanswer, def:physics:phyx-mini-0296:phyxminiproblems-problemphyxmini0296-matchesdisplayedvelocity}
The assigned autoformalize agent should translate this physics problem into Lean declarations in the covered file, with theorem and lemma proof bodies written as `by sorry`.
\end{theorem}
\begin{proof}
This is an autoformalization task, not a proof task. Produce statements that can later be proved by the physics prover without weakening the physical model.
\end{proof}

% --- Archon declaration topology begin ---
\section{Declaration topology}
\begin{definition}[Amplitude Quantity]
  \label{def:physics:phyx-mini-0296:phyxminiproblems-problemphyxmini0296-amplitudequantity}
  \lean{PhyXMiniProblems.ProblemPhyXMini0296.AmplitudeQuantity}
  A nonnegative transverse-displacement amplitude
\end{definition}
\begin{proof}
  This is the declared model definition of Amplitude Quantity.
\end{proof}
\begin{definition}[Wavelength Quantity]
  \label{def:physics:phyx-mini-0296:phyxminiproblems-problemphyxmini0296-wavelengthquantity}
  \lean{PhyXMiniProblems.ProblemPhyXMini0296.WavelengthQuantity}
  A nonnegative standing wavelength
\end{definition}
\begin{proof}
  This is the declared model definition of Wavelength Quantity.
\end{proof}
\begin{definition}[Signed Length Quantity]
  \label{def:physics:phyx-mini-0296:phyxminiproblems-problemphyxmini0296-signedlengthquantity}
  \lean{PhyXMiniProblems.ProblemPhyXMini0296.SignedLengthQuantity}
  A signed axial coordinate or transverse displacement
\end{definition}
\begin{proof}
  This is the declared model definition of Signed Length Quantity.
\end{proof}
\begin{definition}[Time Quantity]
  \label{def:physics:phyx-mini-0296:phyxminiproblems-problemphyxmini0296-timequantity}
  \lean{PhyXMiniProblems.ProblemPhyXMini0296.TimeQuantity}
  A signed time coordinate relative to the graph origin
\end{definition}
\begin{proof}
  This is the declared model definition of Time Quantity.
\end{proof}
\begin{definition}[Duration Quantity]
  \label{def:physics:phyx-mini-0296:phyxminiproblems-problemphyxmini0296-durationquantity}
  \lean{PhyXMiniProblems.ProblemPhyXMini0296.DurationQuantity}
  A nonnegative duration, used here for the oscillation period
\end{definition}
\begin{proof}
  This is the declared model definition of Duration Quantity.
\end{proof}
\begin{definition}[Angular Frequency Quantity]
  \label{def:physics:phyx-mini-0296:phyxminiproblems-problemphyxmini0296-angularfrequencyquantity}
  \lean{PhyXMiniProblems.ProblemPhyXMini0296.AngularFrequencyQuantity}
  Angular frequency, with radians treated as dimensionless
\end{definition}
\begin{proof}
  This is the declared model definition of Angular Frequency Quantity.
\end{proof}
\begin{definition}[Wave Number Quantity]
  \label{def:physics:phyx-mini-0296:phyxminiproblems-problemphyxmini0296-wavenumberquantity}
  \lean{PhyXMiniProblems.ProblemPhyXMini0296.WaveNumberQuantity}
  Wave number, with radians treated as dimensionless
\end{definition}
\begin{proof}
  This is the declared model definition of Wave Number Quantity.
\end{proof}
\begin{definition}[Transverse Velocity Quantity]
  \label{def:physics:phyx-mini-0296:phyxminiproblems-problemphyxmini0296-transversevelocityquantity}
  \lean{PhyXMiniProblems.ProblemPhyXMini0296.TransverseVelocityQuantity}
  A signed transverse-velocity component
\end{definition}
\begin{proof}
  This is the declared model definition of Transverse Velocity Quantity.
\end{proof}
\begin{definition}[nonnegative Length Readout]
  \label{def:physics:phyx-mini-0296:phyxminiproblems-problemphyxmini0296-nonnegativelengthreadout}
  \lean{PhyXMiniProblems.ProblemPhyXMini0296.nonnegativeLengthReadout}
  Read a nonnegative length in the selected length unit
\end{definition}
\begin{proof}
  This is the declared model definition of nonnegative Length Readout.
\end{proof}
\begin{definition}[signed Length Readout]
  \label{def:physics:phyx-mini-0296:phyxminiproblems-problemphyxmini0296-signedlengthreadout}
  \lean{PhyXMiniProblems.ProblemPhyXMini0296.signedLengthReadout}
  \uses{def:physics:phyx-mini-0296:phyxminiproblems-problemphyxmini0296-signedlengthquantity}
  Read a signed axial coordinate or displacement in the selected unit
\end{definition}
\begin{proof}
  This is the declared model definition of signed Length Readout.
\end{proof}
\begin{definition}[time Readout]
  \label{def:physics:phyx-mini-0296:phyxminiproblems-problemphyxmini0296-timereadout}
  \lean{PhyXMiniProblems.ProblemPhyXMini0296.timeReadout}
  \uses{def:physics:phyx-mini-0296:phyxminiproblems-problemphyxmini0296-timequantity}
  Read a signed time coordinate in the selected time unit
\end{definition}
\begin{proof}
  This is the declared model definition of time Readout.
\end{proof}
\begin{definition}[duration Readout]
  \label{def:physics:phyx-mini-0296:phyxminiproblems-problemphyxmini0296-durationreadout}
  \lean{PhyXMiniProblems.ProblemPhyXMini0296.durationReadout}
  \uses{def:physics:phyx-mini-0296:phyxminiproblems-problemphyxmini0296-durationquantity}
  Read a nonnegative duration in the selected time unit
\end{definition}
\begin{proof}
  This is the declared model definition of duration Readout.
\end{proof}
\begin{definition}[angular Frequency Readout]
  \label{def:physics:phyx-mini-0296:phyxminiproblems-problemphyxmini0296-angularfrequencyreadout}
  \lean{PhyXMiniProblems.ProblemPhyXMini0296.angularFrequencyReadout}
  \uses{def:physics:phyx-mini-0296:phyxminiproblems-problemphyxmini0296-angularfrequencyquantity}
  Read angular frequency in radians per selected time unit
\end{definition}
\begin{proof}
  This is the declared model definition of angular Frequency Readout.
\end{proof}
\begin{definition}[wave Number Readout]
  \label{def:physics:phyx-mini-0296:phyxminiproblems-problemphyxmini0296-wavenumberreadout}
  \lean{PhyXMiniProblems.ProblemPhyXMini0296.waveNumberReadout}
  \uses{def:physics:phyx-mini-0296:phyxminiproblems-problemphyxmini0296-wavenumberquantity}
  Read wave number in radians per selected length unit
\end{definition}
\begin{proof}
  This is the declared model definition of wave Number Readout.
\end{proof}
\begin{definition}[transverse Velocity Readout]
  \label{def:physics:phyx-mini-0296:phyxminiproblems-problemphyxmini0296-transversevelocityreadout}
  \lean{PhyXMiniProblems.ProblemPhyXMini0296.transverseVelocityReadout}
  \uses{def:physics:phyx-mini-0296:phyxminiproblems-problemphyxmini0296-transversevelocityquantity}
  Read transverse velocity in coherent selected length/time units
\end{definition}
\begin{proof}
  This is the declared model definition of transverse Velocity Readout.
\end{proof}
\begin{definition}[String Wave Kind]
  \label{def:physics:phyx-mini-0296:phyxminiproblems-problemphyxmini0296-stringwavekind}
  \lean{PhyXMiniProblems.ProblemPhyXMini0296.StringWaveKind}
  The kind of disturbance described in the problem
\end{definition}
\begin{proof}
  This is the declared model definition of String Wave Kind.
\end{proof}
\begin{definition}[Spatial Landmark]
  \label{def:physics:phyx-mini-0296:phyxminiproblems-problemphyxmini0296-spatiallandmark}
  \lean{PhyXMiniProblems.ProblemPhyXMini0296.SpatialLandmark}
  Distinguished axial locations named in the statement or query
\end{definition}
\begin{proof}
  This is the declared model definition of Spatial Landmark.
\end{proof}
\begin{definition}[Standing Wave Role]
  \label{def:physics:phyx-mini-0296:phyxminiproblems-problemphyxmini0296-standingwaverole}
  \lean{PhyXMiniProblems.ProblemPhyXMini0296.StandingWaveRole}
  Standing-wave displacement roles attached to axial landmarks
\end{definition}
\begin{proof}
  This is the declared model definition of Standing Wave Role.
\end{proof}
\begin{definition}[Graph Axis]
  \label{def:physics:phyx-mini-0296:phyxminiproblems-problemphyxmini0296-graphaxis}
  \lean{PhyXMiniProblems.ProblemPhyXMini0296.GraphAxis}
  The two axes of the supplied graph
\end{definition}
\begin{proof}
  This is the declared model definition of Graph Axis.
\end{proof}
\begin{definition}[Axis Quantity]
  \label{def:physics:phyx-mini-0296:phyxminiproblems-problemphyxmini0296-axisquantity}
  \lean{PhyXMiniProblems.ProblemPhyXMini0296.AxisQuantity}
  Physical quantity printed beside each graph axis
\end{definition}
\begin{proof}
  This is the declared model definition of Axis Quantity.
\end{proof}
\begin{definition}[Displacement Graph Landmark]
  \label{def:physics:phyx-mini-0296:phyxminiproblems-problemphyxmini0296-displacementgraphlandmark}
  \lean{PhyXMiniProblems.ProblemPhyXMini0296.DisplacementGraphLandmark}
  Distinguished points of the plotted temporal sinusoid
\end{definition}
\begin{proof}
  This is the declared model definition of Displacement Graph Landmark.
\end{proof}
\begin{definition}[Displacement Time Graph]
  \label{def:physics:phyx-mini-0296:phyxminiproblems-problemphyxmini0296-displacementtimegraph}
  \lean{PhyXMiniProblems.ProblemPhyXMini0296.DisplacementTimeGraph}
  \uses{def:physics:phyx-mini-0296:phyxminiproblems-problemphyxmini0296-amplitudequantity, def:physics:phyx-mini-0296:phyxminiproblems-problemphyxmini0296-signedlengthquantity, def:physics:phyx-mini-0296:phyxminiproblems-problemphyxmini0296-timequantity, def:physics:phyx-mini-0296:phyxminiproblems-problemphyxmini0296-graphaxis, def:physics:phyx-mini-0296:phyxminiproblems-problemphyxmini0296-axisquantity, def:physics:phyx-mini-0296:phyxminiproblems-problemphyxmini0296-displacementgraphlandmark}
  Labels, units, and marked time coordinates supplied by the primary graph. The numerical readings are stated in MatchesPrimaryDisplacementGraph, not built into this structure
\end{definition}
\begin{proof}
  This is the declared model definition of Displacement Time Graph.
\end{proof}
\begin{definition}[Standing String Wave Setup]
  \label{def:physics:phyx-mini-0296:phyxminiproblems-problemphyxmini0296-standingstringwavesetup}
  \lean{PhyXMiniProblems.ProblemPhyXMini0296.StandingStringWaveSetup}
  \uses{def:physics:phyx-mini-0296:phyxminiproblems-problemphyxmini0296-amplitudequantity, def:physics:phyx-mini-0296:phyxminiproblems-problemphyxmini0296-wavelengthquantity, def:physics:phyx-mini-0296:phyxminiproblems-problemphyxmini0296-signedlengthquantity, def:physics:phyx-mini-0296:phyxminiproblems-problemphyxmini0296-timequantity, def:physics:phyx-mini-0296:phyxminiproblems-problemphyxmini0296-durationquantity, def:physics:phyx-mini-0296:phyxminiproblems-problemphyxmini0296-angularfrequencyquantity, def:physics:phyx-mini-0296:phyxminiproblems-problemphyxmini0296-wavenumberquantity, def:physics:phyx-mini-0296:phyxminiproblems-problemphyxmini0296-transversevelocityquantity, def:physics:phyx-mini-0296:phyxminiproblems-problemphyxmini0296-stringwavekind, def:physics:phyx-mini-0296:phyxminiproblems-problemphyxmini0296-spatiallandmark, def:physics:phyx-mini-0296:phyxminiproblems-problemphyxmini0296-standingwaverole, def:physics:phyx-mini-0296:phyxminiproblems-problemphyxmini0296-displacementtimegraph}
  Independent physical quantities and observables for the standing wave. Displacement and transverse velocity remain separate dimensionful fields; the governing-law premise below relates them to the wave parameters
\end{definition}
\begin{proof}
  This is the declared model definition of Standing String Wave Setup.
\end{proof}
\begin{definition}[Matches Standing Wave Problem Data]
  \label{def:physics:phyx-mini-0296:phyxminiproblems-problemphyxmini0296-matchesstandingwaveproblemdata}
  \lean{PhyXMiniProblems.ProblemPhyXMini0296.MatchesStandingWaveProblemData}
  \uses{def:physics:phyx-mini-0296:phyxminiproblems-problemphyxmini0296-signedlengthreadout, def:physics:phyx-mini-0296:phyxminiproblems-problemphyxmini0296-timereadout, def:physics:phyx-mini-0296:phyxminiproblems-problemphyxmini0296-standingstringwavesetup}
  Problem-statement geometry and query data. The graph observes the antinode at x = 0; its adjacent node is at 0.10 m; the requested event is x = 0.20 m, t = 1.0 s. No requested transverse-velocity value occurs in these premises
\end{definition}
\begin{proof}
  This is the declared model definition of Matches Standing Wave Problem Data.
\end{proof}
\begin{definition}[Matches Primary Displacement Graph]
  \label{def:physics:phyx-mini-0296:phyxminiproblems-problemphyxmini0296-matchesprimarydisplacementgraph}
  \lean{PhyXMiniProblems.ProblemPhyXMini0296.MatchesPrimaryDisplacementGraph}
  \uses{def:physics:phyx-mini-0296:phyxminiproblems-problemphyxmini0296-nonnegativelengthreadout, def:physics:phyx-mini-0296:phyxminiproblems-problemphyxmini0296-signedlengthreadout, def:physics:phyx-mini-0296:phyxminiproblems-problemphyxmini0296-timereadout, def:physics:phyx-mini-0296:phyxminiproblems-problemphyxmini0296-durationreadout, def:physics:phyx-mini-0296:phyxminiproblems-problemphyxmini0296-transversevelocityreadout, def:physics:phyx-mini-0296:phyxminiproblems-problemphyxmini0296-standingstringwavesetup}
  Calibrated evidence from the supplied bitmap. Four equal horizontal grid intervals run from 0 s to 2.0 s. The curve passes successively through zero, -y\_s, zero, +y\_s, and zero, where y\_s = 4.0 cm. The slope signs record the visible direction of motion without supplying the queried velocity at x = 0.20 m
\end{definition}
\begin{proof}
  This is the declared model definition of Matches Primary Displacement Graph.
\end{proof}
\begin{definition}[Has Physical Standing Wave Parameters]
  \label{def:physics:phyx-mini-0296:phyxminiproblems-problemphyxmini0296-hasphysicalstandingwaveparameters}
  \lean{PhyXMiniProblems.ProblemPhyXMini0296.HasPhysicalStandingWaveParameters}
  \uses{def:physics:phyx-mini-0296:phyxminiproblems-problemphyxmini0296-nonnegativelengthreadout, def:physics:phyx-mini-0296:phyxminiproblems-problemphyxmini0296-durationreadout, def:physics:phyx-mini-0296:phyxminiproblems-problemphyxmini0296-angularfrequencyreadout, def:physics:phyx-mini-0296:phyxminiproblems-problemphyxmini0296-wavenumberreadout, def:physics:phyx-mini-0296:phyxminiproblems-problemphyxmini0296-standingstringwavesetup}
  Positivity and a principal representative for phase. The phase range does not determine the phase; the graph and governing laws select it
\end{definition}
\begin{proof}
  This is the declared model definition of Has Physical Standing Wave Parameters.
\end{proof}
\begin{definition}[Satisfies Sinusoidal Standing Wave Laws]
  \label{def:physics:phyx-mini-0296:phyxminiproblems-problemphyxmini0296-satisfiessinusoidalstandingwavelaws}
  \lean{PhyXMiniProblems.ProblemPhyXMini0296.SatisfiesSinusoidalStandingWaveLaws}
  \uses{def:physics:phyx-mini-0296:phyxminiproblems-problemphyxmini0296-signedlengthquantity, def:physics:phyx-mini-0296:phyxminiproblems-problemphyxmini0296-timequantity, def:physics:phyx-mini-0296:phyxminiproblems-problemphyxmini0296-nonnegativelengthreadout, def:physics:phyx-mini-0296:phyxminiproblems-problemphyxmini0296-signedlengthreadout, def:physics:phyx-mini-0296:phyxminiproblems-problemphyxmini0296-timereadout, def:physics:phyx-mini-0296:phyxminiproblems-problemphyxmini0296-durationreadout, def:physics:phyx-mini-0296:phyxminiproblems-problemphyxmini0296-angularfrequencyreadout, def:physics:phyx-mini-0296:phyxminiproblems-problemphyxmini0296-wavenumberreadout, def:physics:phyx-mini-0296:phyxminiproblems-problemphyxmini0296-transversevelocityreadout, def:physics:phyx-mini-0296:phyxminiproblems-problemphyxmini0296-standingstringwavesetup}
  Generic laws for a sinusoidal standing wave with an antinode chosen as the spatial origin of phase: k lambda = 2 pi and omega T = 2 pi; the distance from an antinode to its adjacent node is lambda / 4; y(x,t) = y\_m cos(k (x-x\_A)) sin(omega t + phi); transverse particle velocity is the corresponding time derivative. All equations use coherent named-unit readouts. The velocity equation is the dimensionful observable form of Mathlib's scalar sine chain rule.
\end{definition}
\begin{proof}
  This is the declared model definition of Satisfies Sinusoidal Standing Wave Laws.
\end{proof}
\begin{definition}[Answer Choice]
  \label{def:physics:phyx-mini-0296:phyxminiproblems-problemphyxmini0296-answerchoice}
  \lean{PhyXMiniProblems.ProblemPhyXMini0296.AnswerChoice}
  Labels of the four transverse-velocity choices printed with the problem
\end{definition}
\begin{proof}
  This is the declared model definition of Answer Choice.
\end{proof}
\begin{definition}[displayed Velocity In Meters Per Second]
  \label{def:physics:phyx-mini-0296:phyxminiproblems-problemphyxmini0296-displayedvelocityinmeterspersecond}
  \lean{PhyXMiniProblems.ProblemPhyXMini0296.displayedVelocityInMetersPerSecond}
  \uses{def:physics:phyx-mini-0296:phyxminiproblems-problemphyxmini0296-answerchoice}
  Displayed transverse velocity in metres per second
\end{definition}
\begin{proof}
  This is the declared model definition of displayed Velocity In Meters Per Second.
\end{proof}
\begin{definition}[recorded Dataset Answer]
  \label{def:physics:phyx-mini-0296:phyxminiproblems-problemphyxmini0296-recordeddatasetanswer}
  \lean{PhyXMiniProblems.ProblemPhyXMini0296.recordedDatasetAnswer}
  \uses{def:physics:phyx-mini-0296:phyxminiproblems-problemphyxmini0296-answerchoice}
  The answer label recorded in the supplied dataset metadata
\end{definition}
\begin{proof}
  This is the declared model definition of recorded Dataset Answer.
\end{proof}
\begin{definition}[Matches Displayed Velocity]
  \label{def:physics:phyx-mini-0296:phyxminiproblems-problemphyxmini0296-matchesdisplayedvelocity}
  \lean{PhyXMiniProblems.ProblemPhyXMini0296.MatchesDisplayedVelocity}
  \uses{def:physics:phyx-mini-0296:phyxminiproblems-problemphyxmini0296-answerchoice, def:physics:phyx-mini-0296:phyxminiproblems-problemphyxmini0296-displayedvelocityinmeterspersecond}
  Agreement with a velocity displayed to the nearest 0.01 m/s
\end{definition}
\begin{proof}
  This is the declared model definition of Matches Displayed Velocity.
\end{proof}
\begin{lemma}[wavelength In Meters eq two fifths]
  \label{lem:physics:phyx-mini-0296:phyxminiproblems-problemphyxmini0296-wavelengthinmeters-eq-two-fifths}
  \lean{PhyXMiniProblems.ProblemPhyXMini0296.wavelengthInMeters_eq_two_fifths}
  \uses{def:physics:phyx-mini-0296:phyxminiproblems-problemphyxmini0296-nonnegativelengthreadout, def:physics:phyx-mini-0296:phyxminiproblems-problemphyxmini0296-standingstringwavesetup, def:physics:phyx-mini-0296:phyxminiproblems-problemphyxmini0296-matchesstandingwaveproblemdata, def:physics:phyx-mini-0296:phyxminiproblems-problemphyxmini0296-satisfiessinusoidalstandingwavelaws}
  The adjacent-node spacing gives a standing wavelength of 0.40 m
\end{lemma}
\begin{proof}
  The conclusion follows from the governing relations and figure readouts named in the statement of wavelength In Meters eq two fifths.
\end{proof}
\begin{lemma}[wave Number In Radians Per Meter eq five pi]
  \label{lem:physics:phyx-mini-0296:phyxminiproblems-problemphyxmini0296-wavenumberinradianspermeter-eq-five-pi}
  \lean{PhyXMiniProblems.ProblemPhyXMini0296.waveNumberInRadiansPerMeter_eq_five_pi}
  \uses{def:physics:phyx-mini-0296:phyxminiproblems-problemphyxmini0296-wavenumberreadout, def:physics:phyx-mini-0296:phyxminiproblems-problemphyxmini0296-standingstringwavesetup, def:physics:phyx-mini-0296:phyxminiproblems-problemphyxmini0296-matchesstandingwaveproblemdata, def:physics:phyx-mini-0296:phyxminiproblems-problemphyxmini0296-satisfiessinusoidalstandingwavelaws}
  The 0.40 m wavelength gives k = 5 pi rad/m
\end{lemma}
\begin{proof}
  The conclusion follows from the governing relations and figure readouts named in the statement of wave Number In Radians Per Meter eq five pi.
\end{proof}
\begin{lemma}[angular Frequency In Radians Per Second eq pi]
  \label{lem:physics:phyx-mini-0296:phyxminiproblems-problemphyxmini0296-angularfrequencyinradianspersecond-eq-pi}
  \lean{PhyXMiniProblems.ProblemPhyXMini0296.angularFrequencyInRadiansPerSecond_eq_pi}
  \uses{def:physics:phyx-mini-0296:phyxminiproblems-problemphyxmini0296-angularfrequencyreadout, def:physics:phyx-mini-0296:phyxminiproblems-problemphyxmini0296-standingstringwavesetup, def:physics:phyx-mini-0296:phyxminiproblems-problemphyxmini0296-matchesprimarydisplacementgraph, def:physics:phyx-mini-0296:phyxminiproblems-problemphyxmini0296-satisfiessinusoidalstandingwavelaws}
  The graph period 2.0 s gives omega = pi rad/s
\end{lemma}
\begin{proof}
  The conclusion follows from the governing relations and figure readouts named in the statement of angular Frequency In Radians Per Second eq pi.
\end{proof}
\begin{lemma}[phase Offset Radians eq pi]
  \label{lem:physics:phyx-mini-0296:phyxminiproblems-problemphyxmini0296-phaseoffsetradians-eq-pi}
  \lean{PhyXMiniProblems.ProblemPhyXMini0296.phaseOffsetRadians_eq_pi}
  \uses{def:physics:phyx-mini-0296:phyxminiproblems-problemphyxmini0296-standingstringwavesetup, def:physics:phyx-mini-0296:phyxminiproblems-problemphyxmini0296-matchesstandingwaveproblemdata, def:physics:phyx-mini-0296:phyxminiproblems-problemphyxmini0296-matchesprimarydisplacementgraph, def:physics:phyx-mini-0296:phyxminiproblems-problemphyxmini0296-hasphysicalstandingwaveparameters, def:physics:phyx-mini-0296:phyxminiproblems-problemphyxmini0296-satisfiessinusoidalstandingwavelaws}
  At the graph origin the antinode displacement is zero and decreasing. With the principal sine-phase representative, this selects phi = pi
\end{lemma}
\begin{proof}
  The conclusion follows from the governing relations and figure readouts named in the statement of phase Offset Radians eq pi.
\end{proof}
\begin{lemma}[requested Point Spatial Phase eq pi]
  \label{lem:physics:phyx-mini-0296:phyxminiproblems-problemphyxmini0296-requestedpointspatialphase-eq-pi}
  \lean{PhyXMiniProblems.ProblemPhyXMini0296.requestedPointSpatialPhase_eq_pi}
  \uses{def:physics:phyx-mini-0296:phyxminiproblems-problemphyxmini0296-signedlengthreadout, def:physics:phyx-mini-0296:phyxminiproblems-problemphyxmini0296-wavenumberreadout, def:physics:phyx-mini-0296:phyxminiproblems-problemphyxmini0296-standingstringwavesetup, def:physics:phyx-mini-0296:phyxminiproblems-problemphyxmini0296-matchesstandingwaveproblemdata, def:physics:phyx-mini-0296:phyxminiproblems-problemphyxmini0296-satisfiessinusoidalstandingwavelaws}
  The requested point is twice the antinode-to-adjacent-node distance, so its spatial standing-wave phase relative to the graph antinode is pi
\end{lemma}
\begin{proof}
  The conclusion follows from the governing relations and figure readouts named in the statement of requested Point Spatial Phase eq pi.
\end{proof}
\begin{lemma}[transverse Velocity At Requested Event exact]
  \label{lem:physics:phyx-mini-0296:phyxminiproblems-problemphyxmini0296-transversevelocityatrequestedevent-exact}
  \lean{PhyXMiniProblems.ProblemPhyXMini0296.transverseVelocityAtRequestedEvent_exact}
  \uses{def:physics:phyx-mini-0296:phyxminiproblems-problemphyxmini0296-transversevelocityreadout, def:physics:phyx-mini-0296:phyxminiproblems-problemphyxmini0296-standingstringwavesetup, def:physics:phyx-mini-0296:phyxminiproblems-problemphyxmini0296-matchesstandingwaveproblemdata, def:physics:phyx-mini-0296:phyxminiproblems-problemphyxmini0296-matchesprimarydisplacementgraph, def:physics:phyx-mini-0296:phyxminiproblems-problemphyxmini0296-hasphysicalstandingwaveparameters, def:physics:phyx-mini-0296:phyxminiproblems-problemphyxmini0296-satisfiessinusoidalstandingwavelaws}
  At x = 0.20 m, the standing-wave spatial factor is the negative of that at x = 0. At t = 1.0 s the exact model velocity is therefore -0.04 pi m/s = -pi/25 m/s
\end{lemma}
\begin{proof}
  The conclusion follows from the governing relations and figure readouts named in the statement of transverse Velocity At Requested Event exact.
\end{proof}
% --- Archon declaration topology end ---
% --- Archon physics formalization source end ---
